\documentclass{amsart}
\usepackage{amsmath,amssymb,amsthm,hyperref}
\usepackage{algorithm}
\usepackage{algpseudocode}

\newtheorem{theorem}{Theorem}
\newtheorem{lemma}{Lemma}
\newtheorem{proposition}{Proposition}
\newtheorem{corollary}{Corollary}
\theoremstyle{definition}
\newtheorem{definition}{Definition}
\theoremstyle{remark}
\newtheorem{remark}{Remark}

\begin{document}

\title{On the total tardiness of unit jobs under chain precedence:
status of the NP-hardness proof}
\maketitle

\section{The problem}

We are given $n$ unit-length jobs $\{1,\dots,n\}$, all released at time $0$, to be
scheduled nonpreemptively on a single machine. Each job $j$ has an integer due date
$d_j$. A precedence digraph $G$ on $\{1,\dots,n\}$ has maximum indegree $1$ and maximum
outdegree $1$, so $G$ is a disjoint union of directed chains. A feasible schedule is a
permutation $\sigma:\{1,\dots,n\}\to\{1,\dots,n\}$ with $\sigma(j)<\sigma(k)$ for every
arc $(j,k)$ of $G$; here $\sigma(j)$ is the position (equivalently, since every job has
unit processing time and the machine is never idle in any optimal schedule, the
completion time) of job $j$. We minimize the total tardiness
\[
   T(\sigma)\;=\;\sum_{j=1}^{n}\bigl(\sigma(j)-d_j\bigr)^{+},
   \qquad x^{+}:=\max\{0,x\}.
\]
In three--field notation this is $1\mid p_j=1,\ \mathrm{chains}\mid \sum T_j$.

\medskip
\noindent\textbf{Status of this document.}
The intended target is NP-hardness (this is the documented complexity of the problem).
This revision \emph{corrects a fatal error} in the previous write-up: the reduction it
proposed --- both in its original and in its ``corrected'' form (the Remark with item due
dates $N-a_i+r$) --- is \emph{not} a valid reduction. We give below an explicit,
machine-verified counterexample showing the reduction fails, isolate exactly which step is
wrong, and record the parts of the argument that remain correct. We do \emph{not} claim a
complete NP-hardness proof here; the reduction must be replaced by a different gadget, and
the precise nature of the obstruction is documented so that a correct construction can be
built.

\section{What is correct}

\begin{lemma}[Membership in NP]\label{lem:np}
The decision problem (given a threshold $Y$, is there a feasible $\sigma$ with
$T(\sigma)\le Y$?) is in NP.
\end{lemma}

\begin{proof}
A feasible permutation $\sigma$ is a certificate of size $O(n\log n)$. Given $\sigma$ one
checks in time $O(n)$ that $\sigma(j)<\sigma(k)$ for each of the at most $n-1$ arcs of
$G$, and computes $T(\sigma)=\sum_j(\sigma(j)-d_j)^{+}$ in time $O(n)$, comparing it to
$Y$. All numbers occurring are bounded by $n+\max_j|d_j|$, so verification runs in
polynomial time.
\end{proof}

\begin{lemma}[Cost of a contiguous chain block]\label{lem:block}
Let $C$ be a chain $c_1\to\cdots\to c_p$ with staircase due dates
$d_{c_r}=D-p+r$ ($r=1,\dots,p$), for an integer $D$. If, in a schedule, the jobs of $C$
occupy the $p$ consecutive positions $E-p+1,\dots,E$ (so the last job completes at time
$E\ge p$), then $c_r$ is forced into position $E-p+r$ and the tardiness contributed by
$C$ equals
\[
   \sum_{r=1}^{p}\bigl((E-p+r)-(D-p+r)\bigr)^{+}=p\,(E-D)^{+}.
\]
\end{lemma}

\begin{proof}
The arcs $c_r\to c_{r+1}$ force $c_r$ into position $E-p+r$; its tardiness is
$(E-D)^{+}$, independent of $r$, and summing over $r$ gives $p\,(E-D)^{+}$.
\end{proof}

\begin{lemma}[Spreading a staircase block does not help]\label{lem:spread}
If a staircase chain as in Lemma~\ref{lem:block} is placed with its $r$-th job at
position $q_r$ ($q_1<\cdots<q_p$), its tardiness is
$\sum_{r}\bigl(q_r-(D-p+r)\bigr)^{+}$, a coordinatewise nondecreasing function of the
$q_r$. Hence decreasing any $q_r$ (subject to the ordering) never increases the cost, and
a common rightward shift never decreases it.
\end{lemma}

\begin{proof}
Each summand $(q_r-(D-p+r))^{+}$ is nondecreasing in $q_r$.
\end{proof}

\section{The counting argument is also correct}\label{sec:count}

We recall the source problem. \textsc{3-Partition}: given $B$ and $3m$ integers $a_i$
with $B/4<a_i<B/2$ and $\sum_i a_i=mB$, decide whether $\{1,\dots,3m\}$ splits into $m$
triples each summing to $B$. It is strongly NP-complete.

Consider the family of instances used in the previous write-up: set $N=m(B+1)$; for each
$i$ an \emph{item chain} $A_i$ of $a_i$ unit jobs, and one \emph{separator chain}
$P=(p_1\to\cdots\to p_m)$ with $d_{p_t}=t(B+1)$. Suppose a feasible schedule has all
separators on time, i.e.
\begin{equation}\label{eq:deadlines}
   \sigma(p_t)\le t(B+1)\qquad(t=1,\dots,m).
\end{equation}
Let $f(t)$ be the number of item jobs among the first $t(B+1)$ positions. The precedence
$p_1\to\cdots\to p_m$ together with \eqref{eq:deadlines} puts $p_1,\dots,p_t$ all within
the first $t(B+1)$ positions, so $f(t)\le t(B+1)-t=tB$. Symmetrically, the last
$(m-t)(B+1)$ positions hold at most $m-t$ separators, hence at least $(m-t)B$ item jobs,
giving $f(t)\ge mB-(m-t)B=tB$. Thus
\begin{equation}\label{eq:fcount}
   f(t)=tB\quad\text{for all }t,\qquad\text{and}\qquad \sigma(p_t)=t(B+1).
\end{equation}
So \emph{each window} $W_t=\{(t-1)(B+1)+1,\dots,t(B+1)\}$ contains exactly $B$ item jobs
(in its first $B$ slots) and the single separator $p_t$ (in its last slot).
\textbf{This part of the argument is valid.}

\section{The fatal error}\label{sec:error}

The previous write-up then asserted that ``the $B$ item jobs inside a window are exactly a
union of complete item chains,'' i.e.\ that no item chain straddles a window boundary.
\emph{This is false.} A chain may have some jobs in $W_t$ and the rest in $W_{t+1}$, with
the separator $p_t$ sitting (legally) between two of its jobs; feasibility permits this,
and --- crucially --- in the ``corrected'' instance of the Remark, where item due dates
are $d_{a_{i,r}}=N-a_i+r$, item jobs are \emph{never} tardy no matter where they sit. So
straddling carries \emph{zero} cost. Consequently $\eqref{eq:fcount}$ (exactly $B$ item
units per window) is satisfiable by \emph{any} packing, whether or not the $a_i$ admit a
$3$-partition: the partition structure is never used.

\begin{proposition}[Verified counterexample]\label{prop:counter}
Take the strongly-bounded \textsc{3-Partition} instance
\[
   m=2,\quad B=13,\quad (a_1,\dots,a_6)=(4,4,4,4,4,6),
\]
which satisfies $B/4<a_i<B/2$ and $\sum_i a_i=mB=26$, and which is a \emph{NO} instance
(no way to split into two triples each summing to $13$, since five $4$'s and one $6$ admit
no triple summing to $13$). Build the ``corrected'' scheduling instance ($N=m(B+1)=28$;
item due dates $d_{a_{i,r}}=N-a_i+r$; separator due dates $d_{p_t}=t(B+1)$). Then the
schedule
\[
   A_1A_2A_3,\ \underbrace{a_{4,1}}_{\text{pos }13},\ \underbrace{p_1}_{\text{pos }14},\
   \underbrace{a_{4,2}a_{4,3}a_{4,4}}_{15,16,17},\ A_5,\ A_6,\
   \underbrace{p_2}_{\text{pos }28}
\]
(the three jobs $A_1,A_2,A_3$ filling positions $1$--$12$, chain $A_4$ split across the
boundary, $A_5$ at $18$--$21$, $A_6$ at $22$--$27$) is feasible and has total tardiness
$T(\sigma)=0$.
\end{proposition}

\begin{proof}
Every item job sits at a position $\le N=28$ and has due date $N-a_i+r\ge N-a_i+1\ge
28-6+1=23$; in fact each item job's tardiness is $0$ because for the $r$-th job of $A_i$
at position $q$ one checks $q\le N=28$ and $d=N-a_i+r$, and $q-d=q-(N-a_i+r)\le
0$ in this layout (verified directly position by position). The separators are at
positions $14=1\cdot(B+1)$ and $28=2\cdot(B+1)$, exactly their due dates, so they are on
time. Hence $T(\sigma)=0$. (This was confirmed by exhaustive computation: the optimal
total tardiness of this instance is $0$, and likewise for the YES instance
$(4,4,4,4,5,5)$; both are $0$, so the threshold $Y=0$ does not separate YES from NO.)
\end{proof}

\noindent
Proposition~\ref{prop:counter} shows the reduction is invalid: a NO instance of
\textsc{3-Partition} maps to a YES instance of the scheduling decision problem.

\begin{remark}[Why no fix within this gadget family]\label{rem:obstruction}
The obstruction is structural. By \S\ref{sec:count}, once separators meet their
deadlines the item jobs are forced onto the \emph{fixed} position sets
$\{(t-1)(B+1)+1,\dots,(t-1)(B+1)+B\}$, one block of $B$ positions per window. The total
item tardiness then depends only on \emph{which job occupies which of these fixed
positions}, an assignment problem; chain precedence only forbids inverting a chain's
internal order, which never prevents distributing (or splitting) chains across the fixed
slots. With item due dates large enough that items are always on time (needed so that a
\emph{valid} partition gives cost $0$), splitting is free; with item due dates tight
enough to penalize lateness, even a valid partition incurs cost, and that cost depends on
the intra-window ordering of chains rather than on the partition itself. We verified
numerically across the available small NO/YES instances ($B=13$ and $B=20$, $m=2$) that
no choice of item staircase offset $c_0$ (item due $c_0-a_i+r$) combined with separator
due dates $t(B+1)$ yields a threshold cleanly separating YES from NO with a gap that
grows with the instance. Hence the separator-window family of gadgets cannot, by itself,
enforce the $3$-Partition structure; a correct reduction needs a different mechanism that
makes \emph{straddling a boundary} strictly costly (e.g.\ a gadget in which delaying one
job of a chain forces a \emph{later} job, with a tight individual due date, to be tardy).
\end{remark}

\section{Open part}

What remains is to supply a correct polynomial reduction (the target being strong
NP-completeness, matching the documented complexity of
$1\mid p_j=1,\ \mathrm{chains}\mid\sum T_j$). The corrected reduction must avoid the trap
of \S\ref{sec:error}: it cannot rely on separators that merely \emph{count} item units
per window, since chain straddling defeats counting at no cost. A workable gadget should
make a chain's cost \emph{strictly} increase when it is not laid out contiguously inside a
single window --- for instance by giving the chain a tight staircase whose later jobs
become tardy as soon as an earlier job is delayed past a boundary --- so that the
zero-cost (or threshold-cost) schedules are exactly the contiguous, window-aligned ones.
Designing and fully verifying such a gadget is left as the remaining task; the present
document establishes membership in NP (Lemma~\ref{lem:np}), the correct block/spreading
lemmas (Lemmas~\ref{lem:block}--\ref{lem:spread}), the valid window-counting consequence
of separator deadlines (\S\ref{sec:count}), and the precise, machine-verified reason the
proposed reduction fails (Proposition~\ref{prop:counter}, Remark~\ref{rem:obstruction}).

\section{Discussion}

Without precedence constraints, $1\mid p_j=1\mid\sum T_j$ is solved optimally by the
Earliest-Due-Date rule (the cost matrix $c_{j,t}=(t-d_j)^{+}$ is Monge). Chain
constraints make the problem genuinely harder: a direct check shows that the natural
``EDD on ready chain-heads'' greedy is \emph{not} optimal (it fails on small instances),
so the problem is not solved by an obvious priority rule. Establishing NP-hardness
rigorously requires a reduction with a gadget that penalizes chain straddling, which the
previous write-up lacked.

\end{document}
